\section{Background and Related Work}
\label{sec:background}

\subsection{Backpressure Routing}

Tassiulas and Ephremides~\cite{tassiulas1992stability} introduced the \emph{max-weight scheduling}
algorithm for multi-hop wireless networks. At each time slot, the scheduler computes the
differential backlog across each link and activates the schedule that maximizes the
weighted sum of link capacities and backlog differentials. This achieves \emph{throughput
optimality}: any arrival rate vector strictly inside the network capacity region can be
stabilized. Neely~\cite{neely2010stochastic} extended this framework via Lyapunov optimization,
introducing the $V$-parameter tradeoff between queue backlog (delay) and utility maximization.
Georgiadis et al.~\cite{georgiadis2006resource} provide a comprehensive treatment of
cross-layer resource allocation using backpressure principles.

\subsection{Network Pricing and Congestion Control}

Kelly et al.~\cite{kelly1998rate} showed that network rate allocation can be decomposed
into a system problem (shadow prices on link capacities) and user problems (utility
maximization subject to prices). The shadow prices emerge as Lagrange multipliers on
capacity constraints---an insight that directly bridges network routing and economics.
Srikant~\cite{srikant2004mathematics} formalized TCP congestion control as distributed
optimization, with AIMD (Additive Increase Multiplicative Decrease) serving as a
primal-dual algorithm. Jacobson~\cite{jacobson1988congestion} established EWMA for
round-trip time estimation in TCP, providing the practical precedent for our capacity
smoothing mechanism.

\subsection{Token Engineering}

Zhang and Zargham~\cite{zhang2020token} take a dynamical systems approach to token
economy design, modeling token flows as differential equations and using the cadCAD
framework for simulation. Superfluid~\cite{superfluid2024} provides on-chain streaming
payment primitives: Constant Flow Agreements (CFA) for point-to-point streams and
General Distribution Agreements (GDA) for one-to-many proportional splits. BPE adds
a \emph{flow-control layer} above these streaming primitives.

\subsection{Demurrage Economics}

Gesell~\cite{gesell1916natural} proposed demurrage (a holding tax on currency) to
increase monetary velocity. Fisher~\cite{fisher1933stamp} formalized this in \emph{Stamp
Scrip}. Both are \emph{stock-based} velocity mechanisms: they penalize holding money.
BPE is fundamentally different---a \emph{flow-based allocation} mechanism that routes
payments to where capacity exists. The two approaches are orthogonal and composable:
demurrage encourages spending velocity, while BPE ensures efficient allocation of
the resulting flows.

\subsection{Sender-Side Capacity Awareness}

Automated market makers (AMMs) like Uniswap~\cite{adams2021uniswap} implement
sender-side price discovery: the bonding curve encodes available liquidity, and
senders observe the price impact of their trades. This is \emph{sender-side}
capacity awareness. BPE provides \emph{receiver-side} capacity awareness:
sinks (service providers) signal their capacity, and the protocol routes payments
accordingly. The distinction is critical for service economies where the capacity
constraint lies at the receiver, not the liquidity pool.

\subsection{AI Agent Payment Protocols}

The 2024--2026 wave of agent payment protocols addresses authorization and identity:
Google's AP2~\cite{google2025ap2} handles inter-agent authentication,
Coinbase's x402~\cite{coinbase2025x402} uses HTTP 402 for single-shot payments,
OpenAI--Stripe's ACP~\cite{openai2025acp} provides billing and invoicing, and
Visa's TAP~\cite{visa2025tap} extends card-network authorization to agents.
Cha et al.~\cite{cha2025stp} study LLM inference pricing via self-play.
\emph{None of these systems address flow control}---the dynamic rerouting of
payment streams based on real-time capacity. BPE fills this gap.

\subsection{Payment Channel Routing}

Spider~\cite{sivaraman2020spider} applies multi-path transport to payment channel
networks, achieving high throughput by splitting payments across routes and rebalancing
channels. CelerPay~\cite{dong2020celer} generalises state channels with off-chain
token transfers and conditional payment resolution. Malavolta et al.~\cite{malavolta2017concurrency}
address concurrency and privacy in payment-channel networks via hash-time-locked contracts.
These systems optimise \emph{channel-level} routing in bilateral networks. BPE operates
at a different layer: it routes \emph{streaming payment flows} to capacity-constrained
\emph{service providers} via on-chain pool rebalancing. The two approaches are
complementary---BPE could use Spider-style multi-path splitting for the streaming
layer while maintaining backpressure allocation at the service layer.
